\begin{thebibliography}{99}
\interlinepenalty=10000

\bibitem{AljadeffSonnWadsworth} E. Aljadeff, J. Sonn and A. R. Wadsworth,
\emph{Projective Schur groups of Henselian fields},
J. Pure Appl. Algebra 208 (2007), 833--851.
\href{https://doi.org/10.1016/j.jpaa.2006.03.019}{doi:10.1016/j.jpaa.2006.03.019}.

\bibitem{AnnalaHoyoisIwasa} T. Annala, M. Hoyois and R. Iwasa,
\emph{Atiyah duality for motivic spectra}, preprint, 2024,
\href{https://arxiv.org/abs/2403.01561v2}{arXiv:2403.01561}.

\bibitem{BachmannElmantoMorrow} T. Bachmann, E. Elmanto and M. Morrow,
\emph{$\mathbb A^1$-invariant motivic cohomology of schemes}, preprint, 2025,
\href{https://arxiv.org/abs/2508.09915v1}{arXiv:2508.09915}.

\bibitem{BakRehmann1984} A. Bak and U. Rehmann,
\emph{$K_2$-analogs of Hasse's norm theorems},
Comment. Math. Helv. 59 (1984), 1--11.
\href{https://doi.org/10.1007/BF02566335}{doi:10.1007/BF02566335}.

\bibitem{BarnesRational2009} D. Barnes,
\emph{Classifying rational $G$-spectra for finite $G$},
Homology Homotopy Appl. 11 (2009), 141--170.
\href{https://doi.org/10.4310/HHA.2009.v11.n1.a7}{doi:10.4310/HHA.2009.v11.n1.a7}.

\bibitem{BarwickSpectralMackey} C. Barwick,
\emph{Spectral Mackey functors and equivariant algebraic $K$-theory (I)},
Adv. Math. 304 (2017), 646--727.
\href{https://doi.org/10.1016/j.aim.2016.08.043}{doi:10.1016/j.aim.2016.08.043}.

\bibitem{BGS2017} C. Barwick, S. Glasman and J. Shah,
\emph{Spectral Mackey functors and equivariant algebraic $K$-theory (II)},
Tunis. J. Math. 2 (2020), 97--146.
\href{https://doi.org/10.2140/tunis.2020.2.97}{doi:10.2140/tunis.2020.2.97}.

\bibitem{BassTate} H. Bass and J. Tate,
\emph{The Milnor ring of a global field},
in: Algebraic $K$-Theory II: ``Classical'' Algebraic $K$-Theory and Connections with Arithmetic, Lecture Notes in Math. 342, Springer, 1973, pp.~349--446.
\href{https://doi.org/10.1007/BFb0073733}{doi:10.1007/BFb0073733}.

\bibitem{Bernardara2009} M. Bernardara,
\emph{A semiorthogonal decomposition for Brauer--Severi schemes},
Math. Nachr. 282 (2009), 1406--1413.
\href{https://doi.org/10.1002/mana.200610826}{doi:10.1002/mana.200610826}.

\bibitem{BhattScholzeWitt} B. Bhatt and P. Scholze,
\emph{Projectivity of the Witt vector affine Grassmannian},
Invent. Math. 209 (2017), 329--423.
\href{https://doi.org/10.1007/s00222-016-0710-4}{doi:10.1007/s00222-016-0710-4}.

\bibitem{BlochKato1986} S. Bloch and K. Kato,
\emph{$p$-adic \'etale cohomology},
Publ. Math. Inst. Hautes \'Etudes Sci. 63 (1986), 107--152.
\href{https://doi.org/10.1007/BF02831624}{doi:10.1007/BF02831624}.

\bibitem{BlumbergGepnerTabuada} A. J. Blumberg, D. Gepner and G. Tabuada,
\emph{Uniqueness of the multiplicative cyclotomic trace},
Adv. Math. 260 (2014), 191--232.
\href{https://doi.org/10.1016/j.aim.2014.02.004}{doi:10.1016/j.aim.2014.02.004}.

\bibitem{Borel} A. Borel,
\emph{Stable real cohomology of arithmetic groups},
Ann. Sci. Éc. Norm. Supér. (4) 7 (1974), 235--272.
\href{https://doi.org/10.24033/asens.1269}{doi:10.24033/asens.1269}.

\bibitem{BraunChuangLazarev} C. Braun, J. Chuang and A. Lazarev,
\emph{Derived localisation of algebras and modules},
Adv. Math. 328 (2018), 555--622.
\href{https://doi.org/10.1016/j.aim.2018.02.004}{doi:10.1016/j.aim.2018.02.004}.

\bibitem{BHRAdeles} O. Braunling, R. Henrard and A.-C. van Roosmalen,
\emph{A non-commutative analogue of Clausen's view on the id\`ele class group},
J. Inst. Math. Jussieu 23 (2024), 2777--2824.
\href{https://doi.org/10.1017/S1474748024000100}{doi:10.1017/S1474748024000100}.

\bibitem{ChernousovMerkurjev2006} V. Chernousov and A. Merkurjev,
\emph{Connectedness of classes of fields and zero-cycles on projective homogeneous varieties},
Compos. Math. 142 (2006), 1522--1548.
\href{https://doi.org/10.1112/S0010437X06002363}{doi:10.1112/S0010437X06002363}.

\bibitem{ChinburgEtal} T. Chinburg, M. Kolster, G. Pappas and V. Snaith,
\emph{Galois structure of \(K\)-groups of rings of integers},
\(K\)-Theory 14 (1998), 319--369.
\href{https://doi.org/10.1023/A:1007751422988}{doi:10.1023/A:1007751422988}.

\bibitem{Chipchakov2022} I. D. Chipchakov,
\emph{Henselian discrete valued stable fields},
Turkish J. Math. 46 (2022), 1735--1748.
\href{https://doi.org/10.55730/1300-0098.3229}{doi:10.55730/1300-0098.3229}.

\bibitem{ClausenArtin} D. Clausen,
\emph{A $K$-theoretic approach to Artin maps}, preprint, 2017,
\href{https://arxiv.org/abs/1703.07842v2}{arXiv:1703.07842}.

\bibitem{ClausenMathew} D. Clausen and A. Mathew,
\emph{Hyperdescent and étale $K$-theory},
Invent. Math. 225 (2021), 981--1076.
\href{https://doi.org/10.1007/s00222-021-01043-3}{doi:10.1007/s00222-021-01043-3}.

\bibitem{ClausenMathewMorrow} D. Clausen, A. Mathew and M. Morrow,
\emph{$K$-theory and topological cyclic homology of henselian pairs},
J. Amer. Math. Soc. 34 (2021), 411--473.
\href{https://doi.org/10.1090/jams/961}{doi:10.1090/jams/961}.

\bibitem{Deligne1974} P. Deligne,
\emph{La conjecture de Weil. I},
Publ. Math. Inst. Hautes \'Etudes Sci. 43 (1974), 273--307.
\href{https://doi.org/10.1007/BF02684373}{doi:10.1007/BF02684373}.

\bibitem{SGA3} M. Demazure and A. Grothendieck (eds.),
\emph{Sch\'emas en groupes. III: Structure des sch\'emas en groupes r\'eductifs (SGA 3)},
Lecture Notes in Math. 153, Springer, 1970.
\href{https://doi.org/10.1007/BFb0059027}{doi:10.1007/BFb0059027}.

\bibitem{DundasGoodwillieMcCarthy} B. I. Dundas, T. G. Goodwillie and R. McCarthy,
\emph{The Local Structure of Algebraic $K$-Theory},
Algebra and Applications 18, Springer, 2013.
\href{https://doi.org/10.1007/978-1-4471-4393-2}{doi:10.1007/978-1-4471-4393-2}.

\bibitem{DundasMorrow} B. I. Dundas and M. Morrow,
\emph{Finite generation and continuity of topological Hochschild and cyclic homology},
Ann. Sci. Éc. Norm. Supér. (4) 50 (2017), 201--238.
\href{https://doi.org/10.24033/asens.2319}{doi:10.24033/asens.2319}.

\bibitem{GeisserHesselholtNilpotent} T. Geisser and L. Hesselholt,
\emph{On relative and bi-relative algebraic $K$-theory of rings of finite characteristic},
J. Amer. Math. Soc. 24 (2011), 29--49.
\href{https://doi.org/10.1090/S0894-0347-2010-00682-0}{doi:10.1090/S0894-0347-2010-00682-0}.

\bibitem{GeisserLevine} T. Geisser and M. Levine,
\emph{The $K$-theory of fields in characteristic $p$},
Invent. Math. 139 (2000), 459--493.
\href{https://doi.org/10.1007/s002220050014}{doi:10.1007/s002220050014}.

\bibitem{QuillenCurve} D. R. Grayson,
\emph{Finite generation of $K$-groups of a curve over a finite field} (after Daniel Quillen),
in: Algebraic $K$-Theory, Part I (Oberwolfach, 1980), Lecture Notes in Math. 966, Springer, 1982, pp.~69--90.
\href{https://doi.org/10.1007/BFb0062167}{doi:10.1007/BFb0062167}.

\bibitem{Harder} G. Harder,
\emph{Die Kohomologie $S$-arithmetischer Gruppen über Funktionenkörpern},
Invent. Math. 42 (1977), 135--175.
\href{https://doi.org/10.1007/BF01389786}{doi:10.1007/BF01389786}.

\bibitem{HasseNorm1931} H. Hasse,
\emph{Beweis eines Satzes und Widerlegung einer Vermutung über das allgemeine Normenrestsymbol},
Nachr. Ges. Wiss. Göttingen, Math.-Phys. Kl. I (1931), 64--69.

\bibitem{HazratWadsworth} R. Hazrat and A. R. Wadsworth,
\emph{$SK_1$ of graded division algebras},
Israel J. Math. 183 (2011), 117--163.
\href{https://doi.org/10.1007/s11856-011-0045-1}{doi:10.1007/s11856-011-0045-1}.

\bibitem{HesselholtLarsenLindenstrauss2020} L. Hesselholt, M. Larsen and A. Lindenstrauss,
\emph{On the $K$-theory of division algebras over local fields},
Invent. Math. 219 (2020), 281--329.
\href{https://doi.org/10.1007/s00222-019-00909-x}{doi:10.1007/s00222-019-00909-x}.

\bibitem{HopkinsSmith1998} M. J. Hopkins and J. H. Smith,
\emph{Nilpotence and stable homotopy theory II},
Ann. of Math. (2) 148 (1998), 1--49.
\href{https://doi.org/10.2307/120991}{doi:10.2307/120991}.

\bibitem{HuWu2020} Y. Hu and Z. Wu,
\emph{On the Rost divisibility of henselian discrete valuation fields of cohomological dimension 3},
Ann. $K$-Theory 5 (2020), 677--707.
\href{https://doi.org/10.2140/akt.2020.5.677}{doi:10.2140/akt.2020.5.677}.

\bibitem{Izhboldin2000} O. Izhboldin,
\emph{\(p\)-primary part of the Milnor \(K\)-groups and Galois cohomologies of fields of characteristic \(p\)},
in: Invitation to Higher Local Fields, Geom. Topol. Monogr. 3, Geom. Topol. Publ., 2000, pp.~19--29.
\href{https://doi.org/10.2140/gtm.2000.3.19}{doi:10.2140/gtm.2000.3.19}.

\bibitem{Jannsen} U. Jannsen,
\emph{Hasse principles for higher-dimensional fields},
Ann. of Math. (2) 183 (2016), 1--71.
\href{https://doi.org/10.4007/annals.2016.183.1.1}{doi:10.4007/annals.2016.183.1.1}.

\bibitem{KahnPositive} B. Kahn,
\emph{Descente galoisienne et \(K_2\) des corps de nombres},
\(K\)-Theory 7 (1993), 55--100.
\href{https://doi.org/10.1007/BF00962794}{doi:10.1007/BF00962794}.

\bibitem{Kahn2010} B. Kahn,
\emph{Cohomological approaches to $SK_1$ and $SK_2$ of central simple algebras},
Doc. Math. Extra Vol.: Andrei A. Suslin's Sixtieth Birthday (2010), 317--369.
\href{https://doi.org/10.4171/dms/5/10}{doi:10.4171/dms/5/10}.

\bibitem{KahnLevine2010} B. Kahn and M. Levine,
\emph{Motives of Azumaya algebras},
J. Inst. Math. Jussieu 9 (2010), 481--599.
\href{https://doi.org/10.1017/S1474748010000022}{doi:10.1017/S1474748010000022}.

\bibitem{KaroubiWeibelReal} M. Karoubi and C. Weibel,
\emph{Algebraic and real $K$-theory of Real varieties},
Topology 42 (2003), 715--742.
\href{https://doi.org/10.1016/S0040-9383(02)00069-1}{doi:10.1016/S0040-9383(02)00069-1}.

\bibitem{KatoLocalI} K. Kato,
\emph{A generalization of local class field theory by using $K$-groups. I},
J. Fac. Sci. Univ. Tokyo Sect. IA Math. 26 (1979), 303--376.
\href{https://doi.org/10.15083/00077072}{doi:10.15083/00077072}.

\bibitem{KatoLocalII} K. Kato,
\emph{A generalization of local class field theory by using $K$-groups. II},
J. Fac. Sci. Univ. Tokyo Sect. IA Math. 27 (1980), 603--683.
\href{https://doi.org/10.15083/00077073}{doi:10.15083/00077073}.

\bibitem{KatoGaloisCohomology1982} K. Kato,
\emph{Galois cohomology of complete discrete valuation fields},
in: Algebraic $K$-Theory, Part II (Oberwolfach, 1980), Lecture Notes in Math. 967, Springer, 1982, pp.~215--238.
\href{https://doi.org/10.1007/BFb0061904}{doi:10.1007/BFb0061904}.

\bibitem{KatoExistence} K. Kato,
\emph{Existence theorem for higher local fields},
in: Invitation to Higher Local Fields, Geom. Topol. Monogr. 3, Geom. Topol. Publ., 2000, pp.~165--195.
\href{https://doi.org/10.2140/gtm.2000.3.165}{doi:10.2140/gtm.2000.3.165}.

\bibitem{KirsteinKremer} D. Kirstein and C. Kremer,
\emph{A twisted Bass--Heller--Swan decomposition for localising invariants}, preprint, 2024,
\href{https://arxiv.org/abs/2410.22877v1}{arXiv:2410.22877}.

\bibitem{KockZanchettaPower} B. K\"ock and F. Zanchetta,
\emph{Comparison of exterior power operations on higher $K$-theory of schemes},
Math. Z. 309 (2025), Article 63.
\href{https://doi.org/10.1007/s00209-025-03681-2}{doi:10.1007/s00209-025-03681-2}.

\bibitem{Kuku1982} A. O. Kuku,
\emph{Higher algebraic $K$-theory of group-rings and orders in algebras over number fields},
Comm. Algebra 10 (1982), 805--816.
\href{https://doi.org/10.1080/00927878208822751}{doi:10.1080/00927878208822751}.

\bibitem{Kurihara2000} M. Kurihara,
\emph{Kato's higher local class field theory},
in: Invitation to Higher Local Fields, Geom. Topol. Monogr. 3, Geom. Topol. Publ., 2000, pp.~53--60.
\href{https://doi.org/10.2140/gtm.2000.3.53}{doi:10.2140/gtm.2000.3.53}.

\bibitem{LevineHigherChow} M. Levine,
\emph{Bloch's higher Chow groups revisited},
Astérisque 226 (1994), 235--320.

\bibitem{LevineMovingLemma} M. Levine,
\emph{Chow's moving lemma and the homotopy coniveau tower},
$K$-Theory 37 (2006), 129--209.
\href{https://doi.org/10.1007/s10977-006-0004-5}{doi:10.1007/s10977-006-0004-5}.

\bibitem{LevineHomotopyConiveau} M. Levine,
\emph{The homotopy coniveau tower},
J. Topol. 1 (2008), 217--267.
\href{https://doi.org/10.1112/jtopol/jtm004}{doi:10.1112/jtopol/jtm004}.

\bibitem{MathewNaumannNoel2017} A. Mathew, N. Naumann and J. Noel,
\emph{Nilpotence and descent in equivariant stable homotopy theory},
Adv. Math. 305 (2017), 994--1084.
\href{https://doi.org/10.1016/j.aim.2016.09.027}{doi:10.1016/j.aim.2016.09.027}.

\bibitem{Matsumoto} H. Matsumoto,
\emph{Sur les sous-groupes arithmétiques des groupes semi-simples déployés},
Ann. Sci. Éc. Norm. Supér. (4) 2 (1969), 1--62.
\href{https://doi.org/10.24033/asens.1174}{doi:10.24033/asens.1174}.

\bibitem{Merkurjev1995} A. S. Merkurjev,
\emph{$K$-theory of simple algebras},
in: $K$-Theory and Algebraic Geometry: Connections with Quadratic Forms and Division Algebras (Santa Barbara, CA, 1992), Proc. Sympos. Pure Math. 58, Part 1, Amer. Math. Soc., 1995, pp.~65--83.
\href{https://doi.org/10.1090/pspum/058.1/1327281}{doi:10.1090/pspum/058.1/1327281}.

\bibitem{MerkurjevSuslin1982} A. Merkurjev and A. Suslin,
\emph{$K$-cohomology of Severi--Brauer varieties and the norm residue homomorphism},
Izv. Akad. Nauk SSSR Ser. Mat. 46 (1982), 1011--1046, 1135--1136;
English transl., Math. USSR-Izv. 21 (1983), 307--340.
\href{https://doi.org/10.1070/IM1983v021n02ABEH001793}{doi:10.1070/IM1983v021n02ABEH001793}.

\bibitem{MerkurjevSuslinK3} A. Merkurjev and A. Suslin,
\emph{The group $K_3$ for a field},
Izv. Akad. Nauk SSSR Ser. Mat. 54 (1990), 522--545;
English transl., Math. USSR-Izv. 36 (1991), 541--565.
\href{https://doi.org/10.1070/IM1991v036n03ABEH002034}{doi:10.1070/IM1991v036n03ABEH002034}.

\bibitem{MerkurjevSuslin2010} A. Merkurjev and A. Suslin,
\emph{Motivic cohomology of the simplicial motive of a Rost variety},
J. Pure Appl. Algebra 214 (2010), 2017--2026.
\href{https://doi.org/10.1016/j.jpaa.2010.02.006}{doi:10.1016/j.jpaa.2010.02.006}.

\bibitem{MilneEtale} J. S. Milne,
\emph{\'Etale Cohomology},
Princeton Math. Ser. 33, Princeton Univ. Press, 1980.
\href{https://doi.org/10.1515/9781400883981}{doi:10.1515/9781400883981}.

\bibitem{MilneADT} J. S. Milne,
\emph{Arithmetic Duality Theorems},
2nd ed., BookSurge, 2006.

\bibitem{NakayamaMatsushima1943} T. Nakayama and Y. Matsushima,
\emph{\"Uber die multiplikative Gruppe einer $\mathfrak p$-adischen Divisionsalgebra},
Proc. Imp. Acad. Tokyo 19 (1943), 622--628.
\href{https://doi.org/10.3792/pia/1195573246}{doi:10.3792/pia/1195573246}.

\bibitem{NesterenkoSuslin} Yu. P. Nesterenko and A. A. Suslin,
\emph{Homology of the full linear group over a local ring, and Milnor's $K$-theory},
Izv. Akad. Nauk SSSR Ser. Mat. 53 (1989), 121--146;
English transl., Math. USSR-Izv. 34 (1990), 121--145.
\href{https://doi.org/10.1070/IM1990v034n01ABEH000610}{doi:10.1070/IM1990v034n01ABEH000610}.

\bibitem{NikolausScholzeTC} T. Nikolaus and P. Scholze,
\emph{On topological cyclic homology},
Acta Math. 221 (2018), 203--409.
\href{https://doi.org/10.4310/ACTA.2018.v221.n2.a1}{doi:10.4310/ACTA.2018.v221.n2.a1}.

\bibitem{Oka1984} S. Oka,
\emph{Multiplications on the Moore spectrum},
Mem. Fac. Sci. Kyushu Univ. Ser. A 38 (1984), 257--276.
\href{https://doi.org/10.2206/kyushumfs.38.257}{doi:10.2206/kyushumfs.38.257}.

\bibitem{Otabe2023} S. Otabe,
\emph{On the mod $p$ unramified cohomology of varieties having universally trivial Chow group of zero-cycles},
Manuscripta Math. 171 (2023), 215--239.
\href{https://doi.org/10.1007/s00229-022-01381-3}{doi:10.1007/s00229-022-01381-3}.

\bibitem{Platonov1976} V. P. Platonov,
\emph{The Tannaka--Artin problem and reduced $K$-theory},
Izv. Akad. Nauk SSSR Ser. Mat. 40 (1976), 227--261;
English transl., Math. USSR-Izv. 10 (1976), 211--243.
\href{https://doi.org/10.1070/IM1976v010n02ABEH001686}{doi:10.1070/IM1976v010n02ABEH001686}.

\bibitem{QuillenFinite} D. Quillen,
\emph{On the cohomology and $K$-theory of the general linear groups over a finite field},
Ann. of Math. (2) 96 (1972), 552--586.
\href{https://doi.org/10.2307/1970825}{doi:10.2307/1970825}.

\bibitem{QuillenFiniteGeneration} D. Quillen,
\emph{Finite generation of the groups $K_i$ of rings of algebraic integers},
in: Algebraic $K$-Theory I: Higher $K$-Theories, Lecture Notes in Math. 341, Springer, 1973, pp.~179--198.
\href{https://doi.org/10.1007/BFb0067056}{doi:10.1007/BFb0067056}.

\bibitem{QuillenHigher} D. Quillen,
\emph{Higher algebraic $K$-theory: I},
in: Algebraic $K$-Theory I: Higher $K$-Theories, Lecture Notes in Math. 341, Springer, 1973, pp.~85--147.
\href{https://doi.org/10.1007/BFb0067053}{doi:10.1007/BFb0067053}.

\bibitem{ReinerMaximalOrders} I. Reiner,
\emph{Maximal Orders},
London Math. Soc. Monogr. (N.S.) 28, Oxford Univ. Press, 2003.
\href{https://doi.org/10.1093/oso/9780198526735.001.0001}{doi:10.1093/oso/9780198526735.001.0001}.

\bibitem{Richter2006} B. Richter,
\emph{A lower bound for coherences on the Brown--Peterson spectrum},
Algebr. Geom. Topol. 6 (2006), 287--308.
\href{https://doi.org/10.2140/agt.2006.6.287}{doi:10.2140/agt.2006.6.287}.

\bibitem{Riou2010} J. Riou,
\emph{Algebraic $K$-theory, $\mathbf A^1$-homotopy and Riemann--Roch theorems},
J. Topol. 3 (2010), 229--264.
\href{https://doi.org/10.1112/jtopol/jtq005}{doi:10.1112/jtopol/jtq005}.

\bibitem{RognesWeibel} J. Rognes and C. Weibel,
\emph{Two-primary algebraic \(K\)-theory of rings of integers in number fields},
with an appendix by M. Kolster,
J. Amer. Math. Soc. 13 (2000), 1--54.
\href{https://doi.org/10.1090/S0894-0347-99-00317-3}{doi:10.1090/S0894-0347-99-00317-3}.

\bibitem{Rost} M. Rost,
\emph{Chow groups with coefficients},
Doc. Math. 1 (1996), 319--393.
\href{https://doi.org/10.4171/dm/16}{doi:10.4171/dm/16}.

\bibitem{Scholbach2011} J. Scholbach,
\emph{Mixed Artin--Tate motives over number rings},
J. Pure Appl. Algebra 215 (2011), 2106--2118.
\href{https://doi.org/10.1016/j.jpaa.2010.11.019}{doi:10.1016/j.jpaa.2010.11.019}.

\bibitem{SchwedeMoore} S. Schwede,
\emph{Algebraic versus topological triangulated categories},
in: Triangulated Categories, London Math. Soc. Lecture Note Ser. 375, Cambridge Univ. Press, 2010, pp.~389--407.
\href{https://doi.org/10.1017/CBO9781139107075.010}{doi:10.1017/CBO9781139107075.010}.

\bibitem{SerreLocal} J.-P. Serre,
\emph{Local Fields},
Graduate Texts in Math. 67, Springer, 1979.
\href{https://doi.org/10.1007/978-1-4757-5673-9}{doi:10.1007/978-1-4757-5673-9}.

\bibitem{SerreGalois} J.-P. Serre,
\emph{Galois Cohomology},
Springer Monographs in Math., Springer, 2002.
\href{https://doi.org/10.1007/978-3-642-59141-9}{doi:10.1007/978-3-642-59141-9}.

\bibitem{StacksProject} The Stacks Project Authors,
\emph{The Stacks Project},
\href{https://stacks.math.columbia.edu}{https://stacks.math.columbia.edu}.

\bibitem{SuslinClosed} A. Suslin,
\emph{On the $K$-theory of algebraically closed fields},
Invent. Math. 73 (1983), 241--245.
\href{https://doi.org/10.1007/BF01394024}{doi:10.1007/BF01394024}.

\bibitem{Suslin1987} A. Suslin,
\emph{Torsion in $K_2$ of fields},
$K$-Theory 1 (1987), 5--29.
\href{https://doi.org/10.1007/BF00533985}{doi:10.1007/BF00533985}.

\bibitem{SuslinTulenbaev} A. A. Suslin and M. S. Tulenbaev,
\emph{Stabilization theorem for the Milnor $K_2$-functor},
Zap. Nauchn. Sem. Leningrad. Otdel. Mat. Inst. Steklov. (LOMI) 64 (1976), 131--152;
English transl., J. Soviet Math. 17 (1981), 1804--1819.
\href{https://doi.org/10.1007/BF01091768}{doi:10.1007/BF01091768}.

\bibitem{SuslinYufryakov} A. Suslin and A. Yufryakov,
\emph{The $K$-theory of local division algebras},
Dokl. Akad. Nauk SSSR 288 (1986), 832--836;
English transl., Soviet Math. Dokl. 33 (1986), 794--798.

\bibitem{TabuadaVanDenBerghSeparable} G. Tabuada and M. Van den Bergh,
\emph{Noncommutative motives of separable algebras},
Adv. Math. 303 (2016), 1122--1161.
\href{https://doi.org/10.1016/j.aim.2016.08.031}{doi:10.1016/j.aim.2016.08.031}.

\bibitem{TateDuality} J. Tate,
\emph{Duality theorems in Galois cohomology over number fields},
in: Proc. Internat. Congr. Mathematicians (Stockholm, 1962), Inst. Mittag-Leffler, 1963, pp.~288--295.

\bibitem{Voevodsky} V. Voevodsky,
\emph{On motivic cohomology with $\Z/\ell$-coefficients},
Ann. of Math. (2) 174 (2011), 401--438.
\href{https://doi.org/10.4007/annals.2011.174.1.11}{doi:10.4007/annals.2011.174.1.11}.

\bibitem{WadsworthValuation} A. R. Wadsworth,
\emph{Valuation theory on finite dimensional division algebras},
in: Valuation Theory and its Applications I, Fields Inst. Commun. 32, Amer. Math. Soc., 2002, pp.~385--449.
\href{https://doi.org/10.1090/fic/032/21}{doi:10.1090/fic/032/21}.

\bibitem{Wang1950} S. Wang,
\emph{On the commutator group of a simple algebra},
Amer. J. Math. 72 (1950), 323--334.
\href{https://doi.org/10.2307/2372036}{doi:10.2307/2372036}.

\bibitem{WeibelKbook} C. Weibel,
\emph{The $K$-book: An Introduction to Algebraic $K$-theory},
Graduate Studies in Math. 145, Amer. Math. Soc., 2013.
\href{https://doi.org/10.1090/gsm/145}{doi:10.1090/gsm/145}.

\bibitem{Yamazaki2013} T. Yamazaki,
\emph{The Brauer--Manin pairing, class field theory, and motivic homology},
Nagoya Math. J. 210 (2013), 29--58.
\href{https://doi.org/10.1017/S0027763000010722}{doi:10.1017/S0027763000010722}.

\bibitem{Yanchevskii1975} V. I. Yanchevskii,
\emph{The commutator subgroups of simple algebras with surjective reduced norms},
Dokl. Akad. Nauk SSSR 221 (1975), 1056--1058;
English transl., Soviet Math. Dokl. 16 (1975), 492--495.

\bibitem{Zheglov2004} A. B. Zheglov,
\emph{On wild division algebras over fields of power series},
Mat. Sb. 195 (2004), no. 6, 21--56;
English transl., Sb. Math. 195 (2004), 783--817.
\href{https://doi.org/10.1070/SM2004v195n06ABEH000825}{doi:10.1070/SM2004v195n06ABEH000825}.

\end{thebibliography}
